\documentclass[11pt,aps,preprint,nofootinbib,superscriptaddress,tightenlines]{revtex4-2}
\usepackage[T1]{fontenc}
\usepackage{amsmath,amssymb,mathtools}
\let\lambdabar\relax   % revtex defines it; newpxmath does too
\usepackage{newpxtext,newpxmath}
\usepackage{bm}
\usepackage[final]{microtype}
\usepackage{xcolor}
\usepackage{booktabs}
\usepackage{algpseudocode}
\usepackage{dsfont}
\usepackage[most]{tcolorbox}
\newcommand{\one}{\mathds{1}}
\newcommand{\pf}{\bm{\mathsf p}}
\newcommand{\avg}[1]{\langle #1\rangle}
\definecolor{darkblue}{RGB}{20,60,140}
\definecolor{accent}{RGB}{30,80,150}
\definecolor{tint}{RGB}{243,246,251}
% algorithms: a titled, tinted box
\newtcolorbox{algobox}[1]{enhanced, colback=tint, colframe=accent, boxrule=0.6pt, arc=2pt,
  left=6pt, right=6pt, top=4pt, bottom=4pt, fonttitle=\bfseries\small, coltitle=white,
  colbacktitle=accent, title={#1}, before skip=10pt, after skip=10pt}
% the central equations: a tinted panel with an accent rule on the left
\newtcolorbox{keyeq}{enhanced, colback=tint, frame hidden, borderline west={2pt}{0pt}{accent},
  sharp corners, left=4pt, right=4pt, top=0pt, bottom=0pt, before skip=6pt, after skip=6pt}
\usepackage[colorlinks=true,linkcolor=darkblue,citecolor=darkblue,urlcolor=darkblue]{hyperref}

\def\st{\begin{equation}}
\def\stp{\end{equation}}
\def\Eq#1{eq.~(\ref{#1})}
\def\Sect#1{Sect.~\ref{#1}}
\def\App#1{Appendix~\ref{#1}}
\def\tr{\operatorname{tr}}
\let\oldparagraph\paragraph
\renewcommand{\paragraph}{\par\bigskip\oldparagraph}

\begin{document}

\title{The scheme: relativistic stochastic SPH in the density frame}
\author{D.~Teaney}
\affiliation{Department of Physics and Astronomy, Stony Brook University,
Stony Brook, New York 11794, USA}
\date{\today}

\begin{abstract}
A self-contained account of the scheme the code runs by default (\texttt{main}, tag \texttt{defaults-2026-11-15}).  A relativistic fluid is
represented by particles that carry energy, momentum and entropy.  They move by a Hamiltonian step, and they exchange momentum by a
stochastic viscous step built so that its average is the viscous force and its fluctuations are the thermal ones.  Runs start from an exact
equilibrium sampled by hybrid Monte Carlo.  The main text follows that thread; details, special cases and options are in the appendices.
\end{abstract}

\maketitle

%=============================================================================
\section{The idea}
\label{sec:idea}

Fluctuating hydrodynamics has two parts: the ideal equations, which conserve energy and momentum and are reversible, and dissipation,
which produces entropy and comes with thermal noise fixed by the fluctuation--dissipation theorem.  The scheme keeps the two apart and
makes each exact in its own sense.
\begin{itemize}
\item \emph{The ideal step} is a Hamiltonian flow of particles.  It conserves the total energy and momentum and preserves phase-space
  volume, so it leaves the thermal equilibrium of the particles unchanged.
\item \emph{The viscous step} moves momentum between neighbouring particles at fixed positions and particle energies.  It is a Monte Carlo
  chain whose proposals are random momentum transfers.  Their size is set by the viscosity, so that on average they are the viscous
  force; the Metropolis test, with a Hastings factor, makes the thermal equilibrium exactly stationary.
\item \emph{The start} is a sample of that equilibrium, positions included, drawn by hybrid Monte Carlo with the same Hamiltonian.
\end{itemize}
Because each step keeps the same equilibrium, so does the whole: the static fluctuations are right by construction, and the dynamics is
hydrodynamics.

\begin{algobox}{A run}
\begin{algorithmic}[1]
\State draw an equilibrium state: positions, energies and momenta (\Sect{sec:start})
\State viscous step of length $\Delta t/2$
\For{each time step}
  \State ideal step of length $\Delta t$ (\Sect{sec:ideal})
  \State viscous step of length $\Delta t$ (\Sect{sec:viscous}) \Comment{two half steps merged}
  \State record the observables
\EndFor
\end{algorithmic}
\end{algobox}

%=============================================================================
\section{The particles}
\label{sec:particles}

\paragraph{The continuum.}  A relativistic fluid in two space dimensions, written in the density frame: the conserved densities are the energy
density $\varepsilon$ and the momentum density $\bm M$, and the entropy density $S$ is carried along.  The equation of state gives
\st
\varepsilon=h(S,\bm M),\qquad T=\frac{\partial h}{\partial S},\qquad \bm v=\frac{\partial h}{\partial\bm M},\qquad \beta=\frac{\bm v}T .
\stp
$T$ is the bare temperature; the rest-frame temperature is
\st
T_{\rm rel}=\gamma\,T ,
\stp
and appears only with the viscosities.  The equation of state is conformal, and the input is the shear viscosity as $\eta/s$.

\paragraph{The particles.}  Each particle $a$ has a position $\bm x_a$, a momentum $\bm m_a$ and an entropy $s_a$.  Its volume is set by the number
density of particles around it,
\st
V_a=\frac1{n_a},\qquad n_a=\sum_bW(\bm x_a-\bm x_b),
\stp
with $W$ the SPH kernel of size $h$.  The particle's densities are its own entropy and momentum divided by its volume, and its energy is
\st
e_a=V_a\,h\Big(\frac{s_a}{V_a},\frac{\bm m_a}{V_a}\Big),
\stp
with the first law
\st
de_a=T_a\,ds_a+\bm v_a\cdot d\bm m_a-p_a\,dV_a .
\label{firstlaw}
\stp
The entropy per particle sets how much a particle fluctuates: its thermal velocity is of order $1/\sqrt{s_a}$.  The interesting range is
$s_a\approx50$--$500$.

%=============================================================================
\section{The equilibrium}
\label{sec:equilibrium}

At temperature $T$, for a fluid at rest or moving uniformly with velocity $\bm v$, the particles have a definite equilibrium distribution.  It can be
written in two equivalent forms, one in the entropies of the particles and one in their energies.  Each is the natural form for one of the
two steps of the scheme.

\paragraph{In the entropies.}  With the positions, momenta and entropies as variables, and the measure $\prod_ad\bm x_a\,d\bm m_a\,ds_a$,
\begin{keyeq}
\st
\pi\propto\exp\sum_a\Big[s_a-\frac{e_a(s_a,\bm m_a,V_a)-\bm v\cdot\bm m_a}{T}\Big] .
\label{equilibrium}
\stp
\end{keyeq}
The sum of the particle energies is the Hamiltonian $H$ and that of the momenta the total momentum $\bm P$, so the exponent is
$\sum_as_a-(H-\bm v\cdot\bm P)/T$.

\paragraph{In the energies.}  At fixed positions and momenta, change the variable of each particle from its entropy to its energy.  By the first
law \Eq{firstlaw}, at fixed $\bm m_a$ and $V_a$, $ds_a=de_a/T_a$, and the entropy becomes the function $S_a(e_a,\bm m_a,V_a)$.  With the
measure $\prod_ad\bm x_a\,d\bm m_a\,de_a$,
\begin{keyeq}
\st
\pi\propto\exp\sum_a\Big[S_a(e_a,\bm m_a,V_a)-\ln T_a-\frac{e_a-\bm v\cdot\bm m_a}{T}\Big] .
\label{equilibriume}
\stp
\end{keyeq}
The $-\ln T_a$ is the change of measure.  It is not extensive: it is of order one per particle, against an entropy of tens to hundreds.

\paragraph{One form for each step.}
\begin{itemize}
\item \emph{The ideal step} moves the positions and momenta at fixed entropies.  It conserves $H$ and $\bm P$ and preserves $\prod_ad\bm x_a\,d\bm m_a$, so it
  keeps \Eq{equilibrium}.
\item \emph{The viscous step} moves the momenta at fixed positions and energies, and conserves $\bm P$.  On that surface \Eq{equilibriume} is a
  distribution of the momenta alone, \Eq{target}, which the step must keep (\Sect{sec:viscous}).
\end{itemize}
Both steps keep the same distribution, so the run as a whole does.

\paragraph{What it looks like.}  The momenta of the particles fluctuate thermally about the flow, with a variance of order $V_awT$ per component,
and the energies fluctuate with them.  The positions are not a lattice: at finite entropy per particle they form a liquid, with a pair
correlation $g(r)$ and a structure factor that the equilibrium fixes.

\paragraph{What the hydrodynamics must reproduce.}  Since both steps keep this distribution, the equal-time correlations of a run, of the energy
and momentum densities and of the positions, must be those of \Eq{equilibrium} at the same $s_a$, at rest and in a moving fluid.  This is the
test of the scheme's fluctuations: the distribution itself is sampled independently by hybrid Monte Carlo (\Sect{sec:start}), and the
dynamics must agree with it.  It is also how a run starts: from a sample of \Eq{equilibrium}.

%=============================================================================
\section{The ideal step}
\label{sec:ideal}

\paragraph{The Hamiltonian} is the total energy,
\st
H(\bm x,\bm m)=\sum_ae_a ,
\stp
at fixed entropies.  It depends on the positions only through the volumes.  With \Eq{firstlaw}, Hamilton's equations are
\begin{keyeq}
\st
\dot{\bm x}_a=\bm v_a,\qquad
\dot{\bm m}_a=-\sum_b\big(p_aV_a^2+p_bV_b^2\big)\,\bm D_{ab},\qquad
\bm D_{ab}=\frac{\partial W(\bm x_a-\bm x_b)}{\partial\bm x_a} :
\stp
\end{keyeq}
each particle moves with its own velocity, pushed by the symmetric SPH pressure force.  Energy, momentum and angular momentum are
conserved, and the flow preserves the phase-space volume.

\paragraph{The integrator.}  The energy of a particle depends on both its momentum and the positions, so $H$ is not a kinetic plus a potential
energy, and the ordinary leapfrog is then neither symplectic nor second order: the energy drifts.  The generalized leapfrog is symplectic
and time-reversible for any $H$:
\st
\begin{aligned}
\bm m^{1/2}&=\bm m-\frac{\Delta t}2\,\frac{\partial H}{\partial\bm x}\big(\bm x,\bm m^{1/2}\big),\\
\bm x'&=\bm x+\frac{\Delta t}2\Big[\frac{\partial H}{\partial\bm m}\big(\bm x,\bm m^{1/2}\big)+\frac{\partial H}{\partial\bm m}\big(\bm x',\bm m^{1/2}\big)\Big],\\
\bm m'&=\bm m^{1/2}-\frac{\Delta t}2\,\frac{\partial H}{\partial\bm x}\big(\bm x',\bm m^{1/2}\big).
\end{aligned}
\label{glf}
\stp
The first two lines are implicit and are solved by iteration.  The energy no longer drifts; the full step costs about $1.7$ times the
ordinary leapfrog.

%=============================================================================
\section{The viscous step}
\label{sec:viscous}

\subsection{What it must do}

During the viscous step the positions, the volumes and the particle energies are fixed; only the momenta change.  The entropy of a particle
is then a function of its momentum,
\st
S_a(\bm m_a)\equiv S_a(e_a,\bm m_a,V_a),\qquad \frac{\partial S_a}{\partial\bm m_a}=-\beta_a .
\stp
The step must do two things.
\begin{itemize}
\item \emph{On average}, move momentum as viscosity does: the entropy grows at the rate $2K$, with $K$ the dissipative function.
\item \emph{In its fluctuations}, keep the thermal equilibrium.  At fixed energies and positions the equilibrium \Eq{equilibriume} is
\begin{keyeq}
  \st
  \pi(\bm m)\propto\exp\sum_a\big[S_a(\bm m_a)-\ln T_a\big]
  \label{target}
  \stp
\end{keyeq}
  (\Sect{sec:equilibrium}; the flow term drops out because the step conserves the total momentum).
\end{itemize}
A Metropolis chain does both.  A random momentum transfer, drawn from a zero-mean Gaussian with covariance $\Sigma$ and accepted with the
Metropolis probability, drifts on average by
\st
\avg{\Delta\bm m}\simeq\frac12\,\Sigma\,\frac{\partial\ln\pi}{\partial\bm m}=-\frac12\,\Sigma\,\beta .
\stp
The viscous force is $-\Delta t\,\mathcal K\beta$, with $\mathcal K$ the quadratic form of $K$.  The two agree for
\begin{keyeq}
\st
\Sigma=2\Delta t\,\mathcal K :
\stp
\end{keyeq}
the fluctuation--dissipation theorem.  So the noise is read off from $K$, and the Metropolis test keeps the equilibrium.

\subsection{The dissipative function}

\paragraph{Geometry from the kernel, physics from the metric.}  Each particle $a$ has a star, its neighbours within the kernel support.  The
kernel fixes the geometry of the star, from the positions alone: bond weights $w_{ab}$, the bond moment
\st
\mathcal M_a=\sum_{b\neq a}w_{ab}\,\bm r_{ab}\bm r_{ab}^T,\qquad \bm r_{ab}=\bm x_b-\bm x_a ,
\stp
the fitted gradient of $\beta$, and the faces,
\st
(\nabla\beta)_a=\sum_{b\neq a}w_{ab}\,\Delta\beta_{ab}\,\bm r_{ab}^T\mathcal M_a^{-1},\qquad
\bm A_{ab}=V_a\,w_{ab}\,\mathcal M_a^{-1}\bm r_{ab},\qquad
\sum_{b\neq a}\bm A_{ab}\,\bm r_{ab}^T=V_a\one ,
\label{faces}
\stp
with $\Delta\beta_{ab}=\beta_b-\beta_a$.  The last identity, the closure of the faces, holds for any arrangement of the neighbours.  The physics is in
the coefficients: the viscosities and $T_{\rm rel}$, and the comoving metric
\st
\ell_a=\one-\bm v_a\bm v_a^T ,
\stp
which measures the strain in the rest frame of the moving fluid.

\paragraph{The two parts of $K$.}  A linear fit sees the smooth part of the flow; what it misses is the residual of each bond.  $K$ is the sum of
the dissipation of each:
\begin{keyeq}
\st
K=\sum_a\big(K_a^{\rm res}+K_a^{\rm resid}\big).
\stp
\end{keyeq}
The resolved part is the continuum viscous dissipation of the fitted gradient,
\st
K_a^{\rm res}=V_a\,T_{{\rm rel},a}\Big[\eta_a\tr\big(\ell_a\mathring S_a\ell_a\mathring S_a\big)+\tfrac12\zeta_{\rm eff}\,\theta_a^2\Big],
\stp
with $\mathring S_a$ the traceless comoving strain and $\theta_a$ the expansion, both built from $(\nabla\beta)_a$ and $\ell_a$.  The residual part damps the
bond differences left over after the best linear fit,
\st
K_a^{\rm resid}=\frac12\min_X\sum_{b\neq a}c_{ab}\,\big(\Delta\beta_{ab}-X\bm r_{ab}\big)\cdot\ell_a\big(\Delta\beta_{ab}-X\bm r_{ab}\big),
\qquad c_{ab}=2\,V_aT_{{\rm rel},a}\eta_a\,w_{ab}\,f_{ab} ,
\stp
with $f_{ab}$ a factor that gives bonds along the flow the weight of the comoving frame (\App{app:residual}).

\subsection{The two moves}

Each part of $K$ becomes one Monte Carlo move, made by every particle in turn.

\paragraph{The stress move} is the noise of $K^{\rm res}$.  Particle $a$ draws a random symmetric stress $\xi_a$ with covariance $2\Delta t$ times the
Onsager matrix of $K_a^{\rm res}$ (\App{app:hastings}) and gives each bond the traction
\st
\pf_{ab}=\xi_a\bm A_{ab}\qquad(\text{momentum particle }a\text{ gives particle }b),
\stp
giving up their sum itself.

\paragraph{The virial.}  For a set of momentum transfers $\pf_{ab}$ from particle $a$ to its neighbours, the virial is the momentum sent along each
bond times the bond vector, summed over the bonds:
\st
\sum_{b\neq a}\pf_{ab}\,\bm r_{ab}^T .
\stp
It is the momentum flux the transfers carry across the star, $V_a$ times the stress they exert on the fluid, and its antisymmetric part is
their torque.  (\texttt{helfand.tex} writes it with the indices in the other order, $r^i\,\mathsf p^j$; for the symmetric stresses here the two agree.)  For the stress move, the closure \Eq{faces} gives
\st
\sum_{b\neq a}\xi_a\bm A_{ab}\,\bm r_{ab}^T=V_a\,\xi_a :
\stp
the tractions exert exactly the drawn stress, which is symmetric.  The move conserves momentum and angular momentum, and on average it is the
resolved viscous force.

\paragraph{The residual move} is the noise of $K^{\rm resid}$.  Particle $a$ draws an independent Gaussian vector $\mathfrak p_{ab}$ on each bond, with
covariance $2\Delta t\,c_{ab}\,\ell_a$.  Independent draws carry, by chance, a stress of their own, the stress of the draws:
\st
\tilde\xi_a=\frac1{V_a}\sum_{c\neq a}\mathfrak p_{ac}\,\bm r_{ac}^T .
\stp
That stress belongs to the stress move, not to the residual, so it is removed: the draws are conditioned on zero virial.  Each bond gives
back a share of $\tilde\xi_a$ in proportion to its coefficient and its bond vector,
\st
\vec p_{ab}=\mathfrak p_{ab}-c_{ab}\,V_a\tilde\xi_a\,R_a^{-1}\bm r_{ab},\qquad R_a=\sum_{b\neq a}c_{ab}\,\bm r_{ab}\bm r_{ab}^T ,
\stp
with $R_a$ the bond moment weighted by the residual coefficients.  The virial of the residual tractions is then zero for any draw,
\st
\sum_{b\neq a}\vec p_{ab}\,\bm r_{ab}^T=V_a\tilde\xi_a-V_a\tilde\xi_a\,R_a^{-1}R_a=0 :
\stp
the residual tractions exert no stress and no torque.  They move momentum only in patterns that a linear velocity field cannot have,
which the fitted gradient does not see, and on average they are the residual viscous force.

\subsection{The vote}

A proposed move changes the momenta of a star, $\bm m\to\bm m'$, and with them $\sum_a(S_a-\ln T_a)$ by $\Delta\mathbb S$.  The chain keeps the equilibrium
\Eq{target} if every move is balanced by its reverse, detailed balance:
\st
\pi(\bm m)\,q(\bm m\to\bm m')\,P(\bm m\to\bm m')=\pi(\bm m')\,q(\bm m'\to\bm m)\,P(\bm m'\to\bm m),
\label{detbal}
\stp
where $q$ is the probability density of proposing a move and $P$ the probability of accepting it.

\paragraph{Why the Hastings factor: the noise strength depends on the state.}  The covariance of a proposal is set by the coefficients of the
star, $T_{\rm rel}$, $\eta$, $\zeta$ and the metric $\ell$, and these are functions of the momentum of its particle.  The move $\bm m\to\bm m'$ is proposed with
the noise strength of $\bm m$; the move back, $\bm m'\to\bm m$, with that of $\bm m'$.  So the two proposal densities differ,
\st
q(\bm m\to\bm m')\neq q(\bm m'\to\bm m),
\stp
and the plain Metropolis test, which assumes they are equal, violates \Eq{detbal}.  The cure is to put their ratio into the test:
\begin{keyeq}
\st
P=\min\Big(1,\ e^{\Delta\mathbb S}\;\frac{q(\bm m'\to\bm m)}{q(\bm m\to\bm m')}\Big),
\label{vote}
\stp
\end{keyeq}
which satisfies \Eq{detbal} exactly (\App{app:mh}).  The ratio is the Hastings factor.  Without it the error does not shrink with the time step:
the chain becomes a diffusion with a state-dependent strength and no compensating drift, and its equilibrium is wrong by $5$--$20\%$ at
$s_a\approx50$--$15$.  With it, the coefficients are computed at each proposal from the current and from the proposed momentum, and \Eq{target}
is exactly stationary.  Because the faces depend on the positions only, the reverse move always exists and the ratio has a closed form
(\App{app:hastings}).

\begin{algobox}{One viscous step of length $\Delta t$ (positions, volumes and energies fixed)}
\begin{algorithmic}[1]
\State for every particle, its star: bonds, $w_{ab}$, $\mathcal M_a$, faces $\bm A_{ab}$ \Comment{geometry}
\For{the stress sweep and the residual sweep}
  \For{every particle $a$, in a random order} \Comment{parallel: \App{app:sweeps}}
    \State coefficients of $a$ from its current momentum
    \State propose tractions, $\xi_a\bm A_{ab}$ or $\vec p_{ab}$, and the new momenta $\bm m'$ of the star
    \State coefficients at $\bm m'_a$; the Hastings factor; $\Delta\mathbb S$
    \State accept with probability \Eq{vote}
  \EndFor
\EndFor
\State set each particle's entropy to $S_a(e_a,\bm m_a)$
\end{algorithmic}
\end{algobox}

%=============================================================================
\section{The run}
\label{sec:run}

The ideal and the viscous step are combined symmetrically: a viscous half step, an ideal step, a viscous half step.  The two half steps
between consecutive ideal steps are merged into one, as in the box of \Sect{sec:idea}.  The ideal step keeps the equilibrium because it is
Hamiltonian and volume-preserving; the viscous step because it is in detailed balance; so the combination keeps it too.  The time step is
set by the sound speed, the viscosity and the shortest modes of the particles (\App{app:dt}).

%=============================================================================
\section{The start: hybrid Monte Carlo}
\label{sec:start}

A run starts from a sample of the equilibrium, \Eq{equilibrium}, positions included.  Started on a lattice instead, it would relax to the
liquid only slowly, at a rate set by self-diffusion, which falls as $1/s_a$.

\paragraph{Hybrid Monte Carlo} samples it by alternating two exact steps.
\begin{enumerate}
\item \emph{Heat bath} at fixed positions, with the energy form \Eq{equilibriume}.  Each particle draws its energy and momentum from its own
  factor of \Eq{equilibriume}, independently of the others, and its entropy follows, $s_a=S_a(e_a,\bm m_a,V_a)$.
\item \emph{Hamiltonian trajectory} at fixed entropies.  The generalized leapfrog \Eq{glf} runs for a time $\tau$, and the end point is accepted with
  \st
  \min\Big(1,\ e^{-(\Delta H-\bm v\cdot\Delta\bm P)/T}\Big) .
  \stp
  This is the entropy form \Eq{equilibrium}: the exact flow keeps it, and because the integrator is symplectic and reversible, the test removes
  its error exactly.
\end{enumerate}
The trajectories are the fluid's own motion, and they are what moves the positions into the liquid.  A trajectory as long as the time sound
takes to cross the box, $\tau=L/c_s$, decorrelates the longest modes; the length is drawn at random in $[\tau/2,\tau]$, because a fixed length
resonates with a sound mode of the box.

\begin{algobox}{The equilibrium start}
\begin{algorithmic}[1]
\State particles on a hexagonal lattice, with uniform entropy and momentum densities (the flow $\bm v$)
\State exact draw of every particle's energy and momentum at the $T$ and $\bm v$ of the totals
\For{200 updates}
  \State heat bath of every particle at the current positions
  \State trajectory of random length in $[\tau/2,\tau]$, step $0.35$ of the acoustic step; accept or reject
\EndFor
\State exact draw of the energies and momenta at the final positions, shifted so that the total energy and momentum are exact
\end{algorithmic}
\end{algobox}

The trajectories are accepted $95$--$100\%$ of the time.  For $672$ particles the start takes about two minutes.

%=============================================================================
\section{What it achieves}
\label{sec:achieves}

\begin{itemize}
\item \emph{Conservation}: the total momentum exactly, the energy to $10^{-5}$ over $t=100$ at $s_a\gtrsim50$.
\item \emph{Equilibrium}: the equal-time fluctuations of the energy and momentum densities agree with hybrid Monte Carlo of the same model
  to $1$--$2\%$, at rest and in a fluid moving at $v=0.8$.
\item \emph{Dissipation}: shear and sound waves damp at the hydrodynamic rates as $kh\to0$, at rest and moving.
\end{itemize}
The known limits are in \App{app:limits}.

%=============================================================================
\appendix

\section{Detailed balance and the Hastings factor}
\label{app:mh}

\paragraph{Metropolis.}  To sample a distribution $\pi(x)$, propose a move $x\to x'$ with density $q(x\to x')$ and accept it with probability $P(x\to x')$.
The chain keeps $\pi$ if every move is balanced by its reverse,
\st
\pi(x)\,q(x\to x')\,P(x\to x')=\pi(x')\,q(x'\to x)\,P(x'\to x).
\stp
If the proposal is symmetric, $q(x\to x')=q(x'\to x)$, the choice
\st
P(x\to x')=\min\Big(1,\frac{\pi(x')}{\pi(x)}\Big)
\stp
works: of the two moves one is accepted with probability one and the other with $\pi'/\pi$, and both sides equal $\min(\pi,\pi')\,q$.

\paragraph{Hastings.}  If the proposal is not symmetric, the same argument works with the ratio of the proposals included,
\st
P(x\to x')=\min\Big(1,\frac{\pi(x')\,q(x'\to x)}{\pi(x)\,q(x\to x')}\Big):
\stp
both sides of detailed balance are now $\min\big(\pi(x)\,q(x\to x'),\ \pi(x')\,q(x'\to x)\big)$.  The only requirement is that the reverse move can be
proposed at all, $q(x'\to x)>0$.

\paragraph{An example: noise whose strength depends on the state.}  Take a Gaussian proposal of width $\sigma(x)$,
\st
q(x\to x')=\frac1{\sqrt{2\pi}\,\sigma(x)}\exp\Big[-\frac{(x'-x)^2}{2\sigma(x)^2}\Big].
\stp
Forward and back differ through $\sigma$, and the Hastings factor is
\st
\frac{q(x'\to x)}{q(x\to x')}=\frac{\sigma(x)}{\sigma(x')}\exp\Big[-\frac{(x'-x)^2}{2\sigma(x')^2}+\frac{(x'-x)^2}{2\sigma(x)^2}\Big].
\stp
Without it, in the limit of small steps the chain is a diffusion with strength $\sigma(x)^2$ evaluated at the start of each step (It\^o), and its
stationary density is not $\pi$ but
\st
\rho(x)\propto\frac{\pi(x)}{\sigma(x)^2} .
\stp
The bias is set by how much $\sigma$ varies over the distribution, not by the step size, which is why it does not go away with $\Delta t$.  In the
viscous step $x$ is the momentum of a star, $\sigma$ is set by $T_{\rm rel}$, $\eta$ and $\ell$ of its particle, and the many dimensions of the residual
move make the effect larger.

\section{The basis stresses and the Hastings ratios}
\label{app:hastings}

\paragraph{The stress draw.}  With $z_k$ independent standard normal numbers,
\st
\xi_a=\sqrt{\frac{2\Delta t\,T_{{\rm rel},a}}{V_a}}\;\sum_{k=1}^3z_k\,B_k ,
\stp
where $B_1$, $B_2$ are the two traceless symmetric unit matrices dressed with $\ell_a^{1/2}$ on both sides and projected off the scalar, times
$\sqrt{2\eta_a}$, and $B_3\propto\sqrt{\zeta_{\rm eff}}\,\ell_a$.  This gives the covariance $2\Delta t$ times the Onsager matrix of $K_a^{\rm res}$
(\texttt{commonkernel.tex}).

\paragraph{The stress ratio.}  A symmetric $2\times2$ stress has three components; write it as the vector $(\xi^{xx},\xi^{xy},\xi^{yy})$.  Collect the three standard
normal numbers of the draw into a vector
\st
z=(z_1,z_2,z_3),
\stp
and let $J$ be the $3\times3$ matrix whose columns are the scaled basis stresses $\sqrt{2\Delta t\,T_{{\rm rel},a}/V_a}\;B_k$ in these components.  Then
the draw is a linear map of $z$,
\st
\xi_a=J\,z ,
\stp
and the density of proposing $\xi_a$ is that of $z$, a standard Gaussian, divided by the Jacobian $|\det J|$.  The reverse move starts from the
proposed momentum, where the coefficients give $J'$, and must propose the opposite stress, $-\xi_a$.  The numbers it would need are
\st
z'=-J'^{-1}J\,z ,
\stp
and the ratio of the two densities is
\st
\ln\frac{q(\bm m'\to\bm m)}{q(\bm m\to\bm m')}=-\frac12|z'|^2+\frac12|z|^2+\ln|\det J|-\ln|\det J'| .
\stp
If the coefficients did not depend on the momentum, $J'=J$, $z'=-z$, and the ratio would be one.

\paragraph{The residual ratio.}  The density of a Gaussian conditioned on zero virial is the density of the draws divided by that of their
virial at zero.  With $n$ bonds and primes for the proposed state,
\st
\begin{aligned}
\ln\frac{q(\bm m'\to\bm m)}{q(\bm m\to\bm m')}={}&\sum_b\ln\frac{c_{ab}}{c'_{ab}}+\Big(\frac n2-1\Big)\ln\frac{\det\ell_a}{\det\ell'_a}+\ln\frac{\det R'_a}{\det R_a}\\
&-\frac1{4\Delta t}\sum_b\vec p_{ab}\cdot\Big(\frac{\ell_a'^{-1}}{c'_{ab}}-\frac{\ell_a^{-1}}{c_{ab}}\Big)\vec p_{ab} .
\end{aligned}
\stp
Both ratios are checked against direct evaluations of the Gaussian densities in the unit test \texttt{testhastings.jl}.  Each run
records their mean and rms per sweep.

\section{The residual factor and the comoving coefficients}
\label{app:residual}

For a smooth field the residual of a bond is the curvature, weighted by the fourth moment of the star.  With weights round in the lab, bonds
along the flow lack the factor the comoving frame would give them; the residual factor restores it,
\st
f_{ab}=\big(\hat{\bm r}_{ab}\cdot\ell_a\,\hat{\bm r}_{ab}\big)^2 ,
\stp
which is $1/\gamma^4$ for a bond along the flow, $1$ across it, and $1$ at rest.  The viscosities are $\eta_a=(\eta/s)\,s_{\rm rel}$ at the particle's
rest-frame state, and $\zeta_{\rm eff}$ includes the bulk viscosity of the unresolved third dimension.

\section{The sweeps}
\label{app:sweeps}

A move of particle $a$ changes only its star, so moves more than $4h$ apart commute.  The cells of the cell list, at least $2h$ wide, are
coloured so that cells of one colour are at least $4h$ apart.  A sweep visits the colours in random order, the cells of a colour in parallel,
and the particles of a cell in random order, after a random shift of the grid of cells.  The stress and the residual sweep alternate their
order from step to step.  The random numbers of a move come from its particle's own stream, so the result does not depend on the number of
threads.

\section{The time step}
\label{app:dt}

With $\mathrm{CFL}=0.2$, the time step is half the smallest of three limits: sound crossing a kernel, viscous diffusion across a kernel, and the
damping rate $\Gamma_{\max}$ of the shortest modes of the particle representation,
\st
\Delta t=\frac12\min\Big(\mathrm{CFL}\,\frac h{c_s},\ \ \mathrm{CFL}\,\frac{h^2\,T_{\rm rel,min}}{\frac43\,\eta/s},\ \ \frac2{\Gamma_{\max}}\Big).
\stp

\section{The light cone and low entropy per particle}
\label{app:limits}

\paragraph{The cap.}  A state with a Lorentz factor above $\gamma_{\max}=100$ has zero weight; a move or a trajectory that reaches one is rejected.
Near the light cone the inversion of the equation of state loses accuracy as $10^{-16}\gamma^2$, and nothing in a fluid needs $\gamma\gtrsim10$.

\paragraph{Low entropy per particle.}  Below $s_a\approx12$ the thermal motion of a particle is relativistic: many proposals reach the cap, the
residual acceptance falls (to $0.1$ at $s_a=8$), the Hastings log ratios become large, the energy drifts at the $10^{-2}$ level and the shear
damping departs from hydrodynamics.  The equilibrium is still right; the dynamics is no longer hydrodynamic.

\paragraph{Finite $kh$.}  At fixed spacing the shear damping exceeds hydrodynamics by about $0.05\,h/a$ (\texttt{hbya.tex}); whether this is the
renormalization by the thermal fluctuations or a discretization effect is open (\texttt{helfand.tex}).

\section{Options and earlier schemes}
\label{app:options}

Every choice is a setting, recorded in each run's record.
\begin{center}
\begin{tabular}{lll}
\toprule
setting & default & alternatives \\ \midrule
\texttt{SPH\_IDEAL} & \texttt{generalized} & \texttt{particle}, \texttt{grid} \\
\texttt{SPH\_DISSIPATION} & \texttt{moment} & \texttt{common}, \texttt{particle} \\
\texttt{SPH\_HASTINGS} & on where it applies & \texttt{0} \\
\texttt{SPH\_MEASURE} & \texttt{liouville} & \texttt{einstein} \\
\texttt{SPH\_GAMMA\_MAX} & $100$ & \\
\texttt{SR\_POSITIONS} & \texttt{hmc} & \texttt{lattice}, \texttt{poisson} \\
\bottomrule
\end{tabular}
\end{center}
The alternatives:
\begin{itemize}
\item ideal step: \texttt{particle} is the drift--kick--drift leapfrog, \texttt{grid} the step on a fixed lattice;
\item viscous step: \texttt{common} corrects the weights of each star instead of inverting its moment, \texttt{particle} uses weights round in
  the rest frame;
\item measure: \texttt{einstein} is flat in the energies and momenta instead of the entropies; it differs from \Eq{equilibrium} by $\prod_aT_a$, an
  effect of relative size $1/s_a$, and the ideal step does not keep it exactly.
\end{itemize}
Earlier runs are reproduced with
\begin{itemize}
\item 2026-11-14/15: \texttt{SPH\_DISSIPATION=common}, \texttt{SPH\_HASTINGS=0};
\item before 2026-11-14: \texttt{SPH\_DISSIPATION=particle}, \texttt{SPH\_IDEAL=particle},\\
  \texttt{SPH\_MEASURE=einstein}, \texttt{SR\_POSITIONS=lattice}.
\end{itemize}

\section{Further reading}
\label{app:reading}

\begin{itemize}
\item \texttt{energystep.tex}, \texttt{volume\_n.tex}: the Hamiltonian, the particle energy, the number volume.
\item \texttt{equilibrium\_positions.tex}: the equilibrium of the positions, hybrid Monte Carlo, the generalized leapfrog, the measure.
\item \texttt{bond\_dissipation.tex}, \texttt{commonkernel.tex}: the bond form of $K$, the comoving metric, the residual factor.
\item \texttt{stochastic.tex}: the moves, the vote, the sweeps.
\item \texttt{inverse\_moment.tex}: the inverse-moment faces, the conditioned residual, the Hastings correction, and their tests.
\item \texttt{helfand.tex}: the viscosity from Helfand moments and its renormalization.
\end{itemize}

\end{document}
